% Source: https://artofproblemsolving.com/community/c3h2230373_what_is_117
% Stable identity: aops-topic-2230373-q01
\textbf{(a)} Long division gives
\[\dfrac{1}{17}=0.\overline{0588235294117647},\]
a repeating block of length $\boxed{16}$. (Equivalently, $10^{16}-1$ is the first number of the form $10^k-1$ divisible by $17$: $10^8\equiv-1\pmod{17}$.)

\textbf{(b)} In the long division of $1$ by $17$, each step is determined by the current remainder, which is one of $1,2,\ldots,16$ (it is never $0$, since $17\nmid10^k$). After at most $16$ steps a remainder repeats, and from then on the digits repeat. So the period is at most $16$.

\textbf{(c)} Since $100=6\cdot16+4$, the $100$th digit is the $4$th digit of the block $0588\ldots$, namely $\boxed{8}$.
